% Transcription — fonds Alexandre Grothendieck, shelfmark 146, batch 5 (pages 81-100).
% Established from the facsimile archives/batches/146/batch-05.pdf (a mirror of
% https://grothendieck.umontpellier.fr/146.pdf), per the skill
% transcribe-grothendieck.
%
% Pass: Opus 5.5 (claude-opus-5-5), 2026-10-02 — first pass, unchecked against the pages by a human.
%
% Nineteen of the twenty pages are transcribed; page 87 is blank and absent.
%
%   Pages 81-86 — presentations of groupoids: given a morphism of groupoids
%     T~ -> T, generators u~_i lifting generators u_i of T, defining relations
%     M_alpha of T, fibre arrows a_j generating the fibre categories, relations
%     N'_beta among them and conjugation relations (14): claim that the lifted
%     relations (12), (13), (14) define T~. Page 81 is a cancelled first draft
%     of page 82's start.
%   Pages 88-98 — "Généralités sur les groupoïdes": pi-fibrations of groupoids
%     C' -> C, the example of Pi_1(P) -> Pi_1(T) for a U-torsor P, the
%     (pi, mu)-fibrations subordinate to a central extension
%     1 -> pi -> E -> mu -> 1 and the equivalence (17); page 98 is a page of
%     loose notes in two columns.
%   Page 99 — S_0 T_G for an MD graph G with vertices of genus 0 and order <= 4,
%     as a (Z^A, mu_2^A)-groupoid; page 100 a two-line heading only.
%
% The dating is the inventory's own for the shelfmark.
\documentclass[11pt,a4paper]{article}
\input{../preamble/grothendieck}

\folder{146}
\batch{5}
\pages{81}{100}
\dating{[à partir de 1978]}
\watermark{Édition de démonstration}
\foldertitle{Groupoïdes de Teichmüller et les S0 Tg,!ν, S0 TG : notes manuscrites (s.d.).}

\begin{document}

\section{Présentations de groupoïdes}
\note{Titre de l'éditeur. La page 81 ouvre le lot au milieu des notes ; rien n'y renvoie à une page antérieure, mais l'objet (relations définissantes d'un groupoïde) peut prolonger le lot précédent.}

\page{81}
\note{La page est barrée d'un long trait diagonal ; elle se lit néanmoins, et la page 82 reprend le même départ.}

Cas d'un \struck{groupoïde \ill{}} \add{homomorphisme de} groupoïdes
\[
  \widetilde{\mathcal{T}} \xrightarrow{\;p\;} \mathcal{T}
\]
qui est surjectif sur les $\pi_0$. On se donne une famille $(\tilde{u}_i)_{i \in I}$
de flèches de $\widetilde{\mathcal{T}}$, \struck{entre des objets} de projections $u_i$
dans $\mathcal{T}$, et des relations
\note{encadré, avec un renvoi vers $\mathcal{T}$ :} $(i_0'^{(\alpha)}, \ldots, i_{n_\alpha}'^{(\alpha)}) \in \bar{I}'^{\,n_\alpha}$ $(\alpha \in A)$
\[
  (4) \qquad u'_{i'^{(\alpha)}_{n_\alpha}} \cdots u'_{i'^{(\alpha)}_{0}} = 1
  \quad \text{dans } \mathcal{T}
\]
[pour le rôle des relations $R_\alpha$] (flèches composables dans $\mathcal{T}$, avec
$\operatorname{or}(u'_{i'_0}) = \operatorname{ext}(u'_{i'_n})$). De plus, on se donne des
\struck{relations} \add{flèches} $(a_j)_{j \in J}$ de $\widetilde{\mathcal{T}}$ qui sont des
« flèches fibres » (leurs projections dans $\mathcal{T}$ des identités),
\struck{et des relations} d'où des $a'_{j'} = a_j^{\pm 1}$ $(j' \in J \times \{\pm 1\})$, et des
relations
\[
  (5) \qquad a'_{j'^{(\beta)}_{n_\beta}}\, a'_{j'^{(\beta)}_{n_\beta - 1}} \cdots a'_{j'^{(\beta)}_0} = 1
\]
où $j'^{(\beta)} = (j_0'^{(\beta)}, \ldots, j_{n_\beta}'^{(\beta)}) \in J'^{\,n_\beta}$, $\beta \in B$
[pour le rôle des rel. $S_\beta$].

On suppose

1°) Les relations (4) sont une famille de relations définissant $\mathcal{T}$, \struck{en
générateurs} \add{flèches génératrices} $u_i$ $(i \in I)$ et $\operatorname{Ob}\mathcal{T}$ (ce qui
implique que tous les sommets \add{et extr.} des flèches de $\mathcal{T}$ sont parmi les
origines \struck{des \ill{} des (4)} (à fortiori $\widetilde{\mathcal{T}} \to \mathcal{T}$ est
surjectif sur les objets). \struck{et \ill{}}

\marginal{d'où une famille de flèches $\tilde{u}'_{i'} = \tilde{u}_i^{\pm 1}$ paramétrée par
les $i' \in \bar{I}' = I \times \{\pm 1\}$, en posant $\tilde{u}'_{(i,\varepsilon)} =
\tilde{u}_i^{\varepsilon}$ — item pour \ill{}}
\marginal{\textbf{NB} \ill{} d'indices \ill{} $\coprod_{n \in \mathbf{N}^*}$ \ill{} $B$ \ill{} $J'^{\,n}$ \ill{}}

\page{82}
Soit \struck{\ill{}}
\[
  (1) \qquad \widetilde{\mathcal{T}} \longrightarrow \mathcal{T}
\]
un morphisme de groupoïdes, et
\[
  (2) \qquad (\tilde{u}_i)_{i \in I}
\]
une famille de flèches de $\widetilde{\mathcal{T}}$,
\[
  (3) \qquad (u_i)_{i \in I}
\]
la famille correspondante dans $\mathcal{T}$. On suppose que celle-ci est génératrice
(donc $\operatorname{Ob}\mathcal{T}$ est formé des origines et ext. des flèches $u_i$) et on se
donne une famille \struck{de} de relations définissantes de $\mathcal{T}$, en les
flèches génératrices $u_i$. Cela se fait, techniquement, en \uncertain{introduisant}
l'ens. d'indices
\[
  (4) \qquad I' = I \times \{\pm 1\}
\]
et en définissant, pour $i' = (i, \varepsilon) \in I'$
\[
  (5) \qquad u'_{i'} = u_i^{\varepsilon}
\]
et en se donnant une famille $\mathcal{R}_\alpha$, de relations, \uncertain{paramétrée} par un
ens. d'indices $A$, d'éléments de l'ensemble somme
\[
  (6) \qquad \hat{I}' = \coprod_{n \in \mathbf{N}^*} I'^{\,n} .
\]
Donc pour tt $\alpha \in A$, on a un $n_\alpha \in \mathbf{N}^*$,

\page{83}
et un
\[
  (7) \qquad i'^{(\alpha)} \in I'^{\,n_\alpha}, \qquad
  i'^{(\alpha)} = (i'_{\alpha,1}, i'_{\alpha,2}, \ldots, i'_{\alpha,n_\alpha})
\]
d'où une suite d'éléments
\[
  u'_{i'_{\alpha,1}}, u'_{i'_{\alpha,2}}, \ldots, u'_{i'_{\alpha,n_\alpha}} \in \operatorname{Fl}\mathcal{T}
\]
qu'on suppose composables, \struck{\ill{}} et formant un lacet, \struck{\ill{}} et on considère
les relations
\[
  (8) \qquad \underbrace{u'_{i'_{\alpha,1}}\, u'_{i'_{\alpha,2}} \cdots u'_{i'_{\alpha,n_\alpha}}}_{M_\alpha((u_i)_{i \in I})} = 1 .
\]
(ceci posé) pour une telle relation, \struck{considérons les $\tilde{u}'_{i'}$} que j'écris pour
simplifier
\[
  u'_{i'_1}\, u'_{i'_2} \cdots u'_{i'_n} = 1
\]
considérons les flèches \struck{de $\widetilde{\mathcal{T}}$, dans les}
$\tilde{u}'_{i'_1}, \ldots, \tilde{u}'_{i'_n}$ de $\widetilde{\mathcal{T}}$, qui ne sont pas
\uncertain{nécessairement} composables
\[
  (9) \qquad
  \tilde{s}_{i'_n} \to \tilde{t}_{i'_n} \quad
  \tilde{s}_{i'_{n-1}} \to \tilde{t}_{i'_{n-1}} \quad \cdots \quad
  \tilde{s}_{i'_2} \to \tilde{t}_{i'_2} \quad
  \tilde{s}_{i'_1} \to \tilde{t}_{i'_1}
\]
\struck{choisissons, pour tout} Supposons qu'on ait une famille de « flèches »-fibres
\struck{$(a_j)$} \struck{$a_j$}
\[
  (10) \qquad (a_j)_{j \in J}, \qquad a_j \in \operatorname{Fl}(\widetilde{\mathcal{T}})
\]
dans $\widetilde{\mathcal{T}}$, telles que pour tt $s \in \operatorname{Ob}\mathcal{T}$,

\page{84}
les $a_j$ qui sont dans $\widetilde{\mathcal{T}}_s$ forment une famille génératrice de flèches
de $\widetilde{\mathcal{T}}_s$ — ce qui implique que tous les sommets de
$\widetilde{\mathcal{T}}$ sont parmi les \struck{\ill{}} origines et extrémités des $a_j$.
\struck{De l'hypothèse \ill{}} Supposons que $\forall\, \alpha$, \struck{les \ill{} des}
\struck{flèches $\tilde{u}'_{i'_{\alpha,\nu}}$ $(p \leq \nu \leq n_\alpha)$ de
$\widetilde{\mathcal{T}}$,}
$\{\tilde{t}_{i'_n}, \tilde{s}_{i'_{n-1}}\}$, \ldots, $\{\tilde{t}_{i'_2}, \tilde{s}_{i'_1}\}$
dans (9), \struck{qui sont} qui (par \add{suppl. de} composabilité de la suite \add{des}
$u'_{i'}$), sont contenues chacune \uncertain{figurant} dans \uncertain{(9)}, sont contenus
chacun dans un même groupoïde fibre de $\widetilde{\mathcal{T}}$ sur $\mathcal{T}$, soient
contenus dans une même composante connexe de la catégorie fibre correspondante.
L'hyp. sur (10) implique donc qu'il existe dans $\widetilde{\mathcal{T}}$ un mot
\struck{\ill{}} composable en les $\tilde{a}_j$, qui joigne $\tilde{t}_{i'_n}$ à
$\tilde{s}_{i'_{n-1}}$, etc.\ :

\note{Diagramme (11), sur deux lignes ; les flèches étiquetées $M_{\alpha,\ldots}((a_j))$ sont
ondulées sur la page.}

\begin{tikzcd}[column sep=small, row sep=small, nodes={font=\scriptsize}]
\tilde{s}_{i'_n} \arrow[r, "u'_{i'_n}"] & \tilde{t}_{i'_n} \arrow[r, "M_{\alpha,i_n}((a_j))"] & \tilde{s}_{i'_{n-1}} \arrow[r, "u'_{i'_{n-1}}"] & \tilde{t}_{i'_{n-1}} \arrow[r, "M_{\alpha,i_{n-1}}((a_j))"] & \cdots
\end{tikzcd}

\begin{tikzcd}[column sep=small, row sep=small, nodes={font=\scriptsize}]
\cdots & \tilde{s}_{i'_2} \arrow[r] & \tilde{t}_{i'_2} \arrow[r, "M_{\alpha,2}((a_j))"] & \tilde{s}_{i'_1} \arrow[r, "u'_{i'_1}"] & \tilde{t}_{i'_1}
\end{tikzcd}

\page{85}
La flèche composée de (11) est alors une flèche
$\tilde{s}_{i'_n} \to \tilde{t}_{i'_1}$ qui se trouve dans \struck{\ill{}} une catégorie fibre,
\add{la flèche \uncertain{apparaissant} \ill{}} \ill{} trouvons un mot \add{$M_{\alpha,i_1}$} en
les $a_j$ tels que \struck{$M$}
\[
  M_{\alpha,1}((a_j)) : \tilde{t}_{i'_1} \rightsquigarrow \tilde{s}_{i'_n}
\]
et que le composé de \struck{la flèche} celle-ci avec la \uncertain{précédente} soit
l'identité :
\[
  (12) \qquad
  \boxed{\begin{gathered}
    M_{\alpha,1}((a_j))\, u'_{i'_{\alpha,1}}\, M_{\alpha,2}((a_j))\, u'_{i'_{\alpha,2}} \cdots \\
    \cdots M_{\alpha,i'_{\alpha,n_\alpha}}((a_j))\, u'_{i'_{\alpha,n_\alpha}} = 1
  \end{gathered}}
\]
\marginal{relation (8) « remontée »}

Supposons donné de plus une \struck{\ill{}} \add{famille} de relations sur les flèches-fibres
$a_j$
\[
  (13) \qquad \boxed{N'_\beta((a_j)) = 1}
\]
telles que celles de ces relations qui sont relations dans une catégorie-fibre fixée
$\widetilde{\mathcal{T}}_s$, soient définissantes pour celle-ci \struck{famille génératrice}
(compte tenu \struck{\ill{}} de flèches \struck{\ill{}} $(a_j)$). \struck{Supposons} Supposons
enfin que pour tt $\tilde{s} \in \widetilde{S}$, les \struck{flèches} flèches $a_j$ qui sont
des lacets en $\tilde{s}$ forment une

\page{86}
syst. de générateurs du \add{groupe $\operatorname{Aut}_{\widetilde{\mathcal{T}}_s}(\tilde{s})$, où
$\widetilde{\mathcal{T}}_s$ est} la catégorie fibre, et donnons-nous de plus, pour tt
$a_{j_0}$ qui est un lacet, et tout $u_i$ qui a même origine $\tilde{s}$ que $a_{j_0}$, une
relation
\[
  (14) \qquad
  \boxed{\underbrace{\tilde{u}_i(a_{j_0})}_{\substack{\text{flèche de la}\\ \text{catégorie fibre}\\ \text{contenant ext.}(\tilde{u}_i)}}
  = M''_{i,j}((a_j))} .
\]
Cela \struck{\ill{}} posé, je dis que (12) (13) (14) est un syst. de relations définissant
pour $\mathcal{T}$.
\note{On attendrait $\widetilde{\mathcal{T}}$ ; la page porte le $\mathcal{T}$ sans tilde. Le reste de la page est blanc, et l'argument s'arrête là : la page 87 est blanche, et la page 88 ouvre un autre développement.}

\section{Généralités sur les groupoïdes}
\note{Titre de sa main, souligné, en tête de la page 88.}

\page{88}
Soient \struck{\ill{}} \add{$C, C'$} des groupoïdes, et
\[
  (1) \qquad f : C' \longrightarrow C
\]
un morphisme, enfin $\pi$ un groupe \struck{abélien} \add{(\uncertain{non nécessairement}
abélien)}. On dit que $f$ est une \add{cofibr. surjective} $\pi$-fibration \add{(si}
\add{\uncertain{bijective})}

a) $f$ est \add{ent} surjection \add{sur les objets, et (sur les flèches)}
\add{\uncertain{bijective} \ill{}} \ill{}

b) Les fibres de $f$ (qui sont des groupoïdes connexes) sont munies d'une structure de
$\pi$-groupoïdes \uncertain{connexes} \ill{}, i.e. le groupe fondamental est commutatif, et on
donne un isomorphisme entre celui-ci et $\pi$.

Donc on a un homomorphisme de groupes
\[
  (2) \qquad \pi \longrightarrow \operatorname{End}(\mathrm{id}_{C'})
\]
i.e.\ $\pi$ opère fonctoriellement sur tt objet de $C'$, et on a une suite exacte
\[
  1 \longrightarrow \pi \longrightarrow \operatorname{Aut}_{C'}(x) \longrightarrow
  \operatorname{Aut}_{C'}(f(x)) \longrightarrow 1 .
\]
\note{Le dernier groupe est écrit $\operatorname{Aut}_{C'}$ sur la page ; on attendrait
$\operatorname{Aut}_C$.}

(3) Les données (1), (2), satisfaisant : la \struck{surjectivité de $f$ sur les objets} \ill{}
\struck{bijectivité de $\pi_0(f)$ \ill{} $\pi_0$ \ill{}}, \ill{} plus la bijectivité de
$\pi_0(f)$ \ill{} — \ill{} strictes si de plus $f$ est surjective sur les objets.
\[
  \pi_0(f) : \pi_0(C) \xrightarrow{\;\sim\;} \pi_0(C') .
\]

(4) \textbf{Ex} Soient $U$ un groupe top. \struck{\ill{}} connexe et loc.\ 1-connexe, $P$ un
$U_T$-torseur sur un espace $T$,

\marginal{qui fait de $C'$ une \emph{extension} de $C$ par $\pi$}
\marginal{l'extension est ici néc.\ centrale ; si on veut des cat.\ (3) plus générales,
\uncertain{non nécessairement} centrales, il faut remplacer le groupe constant $\pi$ par un
syst.\ local de gr.\ sur $C$ \ill{} on peut \ill{}}

\page{89}
alors pour les comp.\ connexes $T_i$ de $T$,
\[
  \text{\struck{$\pi_1(P_i, x_i) \to$}}
\]
\[
  \text{\struck{\ill{}}}\quad \pi \longrightarrow \pi_1(P_i, x_i) \longrightarrow \pi_1(T_i, t_i) \longrightarrow 1
\]
où $\pi = \pi_1(U, 1)$, \uncertain{supposons} que les \uncertain{homomorphismes}
$\pi \to \pi_1(P_i, x_i)$ soient injectifs, alors
\[
  \Pi_1(p) : \Pi_1(P) \longrightarrow \Pi_1(T)
\]
\uncertain{est} muni d'une $\pi$-structure. (Si $\pi \to \pi_1(P_i, x_i)$ \struck{était} n'est
pas supposé injectif, \add{et} que son noyau $N_i$ est indép.\ de $i$, $\Pi_1(p)$ est
\struck{\ill{}} une $\pi'$-fibration, $\pi' = \pi / N_i$.)

Soit, dans l'exemple précédent, $E$ un ensemble \struck{commutatif} sur qui opère un
groupe \struck{\ill{}} $\mu$,
\[
  (6) \qquad \mu \xrightarrow{\;i_\mu\;} U
\]
un hom.\ de groupes, et
\[
  (7) \qquad E \xrightarrow{\;i_E\;} P
\]
une application, compatible avec les opérations de $\mu$. Si on regarde le \struck{groupoïde}
groupoïde $C'$ sur $E$ comme ens.\ d'objets, induit par $\Pi_1(P)$ et (7),

\marginal{$x_i \in P_i = P | T_i$, $t_i$ image dans $T_i$ (5)}
\marginal{$p : P \to T$ surjectif \ill{} à fibres \uncertain{connexes} \ill{}}
\marginal{\textbf{NB} Si on suppose $U$ \uncertain{connexe} \ill{} $\pi_1(U, 1)$ \ill{}
$\pi_1(U_T)$ \ill{} \uncertain{localement constant} \ill{} $\pi_1(P)$ \ill{}
$\pi_0$-structure \ill{}}
\note{L'ensemble $E$ est tracé d'une lettre appuyée, distincte du $\mathcal{E}$ cursif qu'il emploie à partir de la page 91 pour l'extension (13) ; la transcription garde la distinction.}
\marginal{\textbf{NB} comme $\mu$ opère librement sur $P$, il faut que $i_\mu$ soit
\uncertain{injectif}, et que $\mu$ opère librement sur $S$.}
\note{Les notes de la marge gauche sont écrites en biais et en partie recouvertes ; seuls les
fragments lisibles sont donnés.}

\page{90}
\[
  (8) \qquad \operatorname{Hom}_{C'}(x, y) = \operatorname{Hom}_{\Pi_1(P)}(i_E(x), i_E(y))
  \qquad x, y \in E
\]
On voit par transport de structure que $\mu$ opère sur $C'$, de façon que
\[
  C' \longrightarrow \Pi_1(P)
\]
soit compatible avec les op.\ de $\mu$.

Considérons l'application
\[
  (9) \qquad E_0 = E/\mu \xrightarrow{\;i_{E_0}\;} T = P/U
\]
induite par passage aux quotients, et soit \struck{\ill{}} $C$ le groupoïde sur $E_0$ comme
ens.\ d'objets induit par $i_{E_0}$ :
\[
  (10) \qquad \operatorname{Hom}_C(x_0, y_0) = \operatorname{Hom}_{\Pi_1(T)}(i_{E_0}(x_0), i_{E_0}(y_0)) .
\]
On trouve alors un diagramme commutatif d'hom.\ de groupoïdes
\note{Diagramme (11). Sur la flèche horizontale du haut et à côté de $f$ figurent deux mentions
biffées, illisibles.}

\begin{tikzcd}
C' \arrow[r] \arrow[d, "f"'] & \Pi_1(P) \arrow[d] \\
C \arrow[r] & \Pi_1(T)
\end{tikzcd}

compatible avec les op.\ de $\mu$, en faisant

\marginal{les flèches horiz.\ sont des équivalences, si $E \to \pi_0(T)$ surjectif \ldots}

\page{91}
opérer $\mu$ trivialement sur $C$ et sur $\Pi_1(T)$.

On en déduit les op.\ de $\mu$ sur $\Pi_1(P)$ resp.\ sur $C'$ avec la $\pi$-structure sur
$\Pi_1(p)$, \uncertain{celle} de $f$ qui s'en déduit par transport de structure. Considérons
pour ceci la suite exacte
\[
  (12) \qquad 0 \longrightarrow \pi \longrightarrow \widetilde{U} \longrightarrow U \longrightarrow 0
\]
et son image inverse par (6), d'où
\[
  (13) \qquad 0 \longrightarrow \pi \longrightarrow \mathcal{E} \longrightarrow \mu \longrightarrow 0
  \qquad \text{ext.\ \uncertain{commutative}}
\]
\struck{On suppose} \struck{Supposons que $\mu$}
Soient $x, y$ deux points de $E$ \struck{\ill{}} qui \add{i.e.\ dans} \add{une même fibre
$C'_\xi$ de $C'$} \ill{} soient \emph{congrus mod} $\mu$, (\struck{donc} on suppose que
\struck{\ill{}} $\mu$ opère librement sur $E$ \uncertain{(cela équivaut à $i_E$ injectif)}.
Je dis qu'on a alors un isom.\ canonique de $\pi$-torseurs
\[
  (14) \qquad \operatorname{Isom}_{C'_\xi}(x, y) \simeq \mathcal{E}_{y - x}
\]
où $y - x$ est l'unique $z \in \mu$ tel que $y = z \cdot x$, et où $\mathcal{E}_z$ désigne la
fibre de l'ext.\ $\mathcal{E}$ (13) sur l'élément $z \in \mu$. Cet isom.\ (14) est compatible
avec la composition dans $C'_\xi$, qui correspond à l'addition dans $\mathcal{E}$.

\page{92}
En d'autres termes, on trouve un isom.\ \add{de $\pi$-groupoïdes}
\[
  (15) \qquad C'_\xi \simeq \mathcal{E}(E_\xi)
\]
avec les notations introduites plus haut, où $E_\xi$ est le $\mu$-torseur image inverse de
$\xi \in E_0 = \operatorname{Ob} C$ par $E \to E_0 = E/\mu$.

De façon générale, étant donné une extension \add{centrale} de groupes
\struck{commutatif} $\mathcal{E}$ (13), une \struck{\ill{}} $(\pi, \mu)$-fibration
\add{(stricte)} de groupoïdes \emph{subordonnée} : l'extension (13), \uncertain{consiste en}
\uncertain{un} hom.\ de groupoïdes
\[
  (16) \qquad f : C' \longrightarrow C
\]
\struck{muni}

2°) D'une op.\ \add{(externe)} de $\mu$ sur $C'$, \struck{de façon} \add{compatible avec les
$\pi$-opérations internes}

1°) D'une op.\ (internes) de $\pi$ sur $C'$, i.e.\ d'un hom
\[
  \pi \longrightarrow \text{\struck{$\operatorname{Aut}$}}\,(\mathrm{id}_{C'})
\]
satisfaisant les conditions suivantes
\note{Une accolade réunit 2°) et 1°) et les renvoie par une flèche à « satisfaisant les conditions suivantes ».}

a) La structure 1°) fait de $f$ une \add{$\pi$-}fibration stricte, en particulier \ill{}
\uncertain{contenue} \ill{} (13), et bijectivité de $\pi_0(f)$ [$+$ surjectivité de $f$ sur
les objets, qu'il faut préciser]

\page{93}
b) \struck{\ill{}} $f$ est compatible avec opérations de $\mu$, quand on fait opérer $\mu$
trivialement sur $C$ \struck{\ill{}}.

c) $\mu$ opère librement sur $\operatorname{Ob} C'$, et
$\operatorname{Ob} C'/\mu \to \operatorname{Ob} C$ est bijectif ($\Rightarrow f$ surjectif sur
les objets)

d) \struck{Exactement} Condition à examiner, sur la structure des catégories fibres
$C'_\xi$ $(\xi \in \operatorname{Ob} C)$, en tant que $\pi$-\add{connexes} groupoïdes
\struck{\ill{}}, avec groupe d'opérateurs $\mu$.

Soit $\Phi$ \struck{est} un $\pi$-groupoïde \add{0-connexe}, avec op.\ de $\mu$ sur
\uncertain{$C'$} (compatible avec $\pi$ structure), tel que $\operatorname{Ob}\Phi$ soit un
torseur sous $\mu$, soit $E$. Soient $x, y \in E$, et considérons
$\operatorname{Hom}_\Phi(x, y)$, c'est un $\pi$-torseur, soit $H_{x,y}$, et on a, via l'op.\
de $\mu$
\[
  H_{x,y} \xrightarrow[\;\sim\;]{\;\alpha_*\;} H_{\alpha x, \alpha y} \qquad \alpha \in \mu
\]
On trouve un système transitif d'isom.\ entre les $H_{x,y}$ pour lesquels
$y^{-1} x = u \in \mu$ est fixé ; donc par là identifiés entre eux à un $H_u$.
\struck{\ill{}} La composition des homomorphismes s'exprime par des isom.\ canoniques
\[
  H_u \wedge H_v \xrightarrow{\;\sim\;} H_{uv}
\]
satisfaisant la condition d'associativité \add{d'unité}, qu'il

\page{94}
faut pour définir une extension
\[
  1 \longrightarrow \pi \longrightarrow
  \underbrace{\mathcal{E}_\Phi}_{\substack{\Phi \text{ avec sa}\\ (\pi,\mu)\text{ structure}}}
  \longrightarrow \mu \longrightarrow 1
\]
\uncertain{CQFD} (on n'a pas utilisé la commutativité de $\mu$ encore) — et \add{inversement,}
\struck{une telle} \add{une} extension (\uncertain{pas} \uncertain{nécessairement} commutative)
\[
  1 \longrightarrow \pi \longrightarrow \mathcal{E} \longrightarrow \mu \longrightarrow 1
\]
\uncertain{plus} la donnée d'un torseur $E$ sous $\mu$, permet de construire canoniquement un
$\pi$-groupoïde $\Phi = \mathcal{E}(E)$, plus (par transport de structure) une opération de
$\mu$ sur $\Phi = \mathcal{E}(E)$. En résumé, on a \struck{une} une équivalence de catégories
\[
  (17) \qquad
  \left[\begin{array}{l}
    \pi\text{-groupoïdes conn.\ } \Phi\text{, avec}\\
    \text{op.\ du groupe } \mu \text{ sur}\\
    (\Phi, \pi)\text{, tel que } \operatorname{Ob}\Phi\\
    \text{soit un } \mu\text{-torseur}
  \end{array}\right]
  \xrightarrow{\;\approx\;}
  \left[\begin{array}{l}
    \text{extensions de}\\
    \mu \text{ par } \pi
  \end{array}\right]
\]
\note{« conn.\ $\Phi$ » est ajouté par lui au-dessus de la ligne. Les crochets sont de l'éditeur ; sur la page les deux catégories sont écrites en colonnes de part et d'autre de la flèche.}

Ceci dit, l'extension (13) étant donnée une fois pour toutes, on dit qu'on a une
$\mathcal{E}(\pi, \mu)$ structure sur $C'$ au-dessus de $C$, si \struck{\ill{}} \add{en} plus
des conditions a) b) c)) pour tout $\xi \in \operatorname{Ob} C$ on se donne un iso.\ entre

\marginal{\uncertain{Pour} les hyp.\ faites, cette extension est \uncertain{centrale},
\uncertain{mais} \ill{} pour \uncertain{opérer} $\mu$ \uncertain{non trivialement} sur $\pi$
\ldots}
\marginal{la commutativité de $\pi$ \ill{} torseurs \ill{} simplifie \ill{}}

\page{95}
l'extension $\mathcal{E}_{C'_\xi}$ de \struck{$\mu$} \add{$\mu$} par $\pi$, et l'extension
$\mathcal{E}$, avec une condition de passage d'un \struck{\ill{}} $\xi$ à un autre, je présume,
pour $\xi_1, \xi_2 \in \operatorname{Ob} C$ dans une même composante connexe de $C$.

Pour préciser ce point, je vais m'inspirer de l'exemple précédent avec $C'$, $C$ décrits en
termes de $P/T$ et des \struck{\ill{}}homomorphismes
\[
  E \longrightarrow P, \qquad \mu \longrightarrow U,
\]
en notant que l'on a alors un diagramme commutatif

\begin{tikzcd}
E \arrow[r] \arrow[d] & P \arrow[d, "q"] \\
E/\mu = E_0 \arrow[r] & P_0 = P/\mu
\end{tikzcd}

\note{Après $P_0 = P/\mu$, une mention biffée, illisible.}
Notons que \struck{pour} $P_0$ est un $U_0$-torseur sur $T$, avec
\[
  (19) \qquad U_0 = U/\mu
\]
et par suite on a $\widetilde{U}_0 = \widetilde{U}$ et un iso can.
\[
  (20) \qquad \mathcal{E} \simeq \pi_1(U_0, 1_{U_0}) .
\]
\note{Le $\mathcal{E}$ est écrit au-dessus d'un $\Pi$ biffé.}
Par suite, \struck{\ill{}} $\Pi_1(P_0)$ est muni d'une

\marginal{On suppose $\mu \hookrightarrow U$ (18), et $\mu$ distingué dans $U$}
\marginal{\textbf{NB} \uncertain{On a donc} $C$ : \ill{} quotient \ill{}
$C \simeq C'/\mathcal{E}(\pi, \mu)$ \ill{}}

\page{96}
$\mathcal{E}$-structure
\[
  (21) \qquad \mathcal{E} \longrightarrow \operatorname{Aut} \mathrm{id}_{\Pi_1(P_0)}
\]
d'où si $x_0, y_0 \in P_0$, $\mathcal{E}$ opère librement sur
$\operatorname{Hom}_{\Pi_1(P_0)}(x_0, y_0)$ à droite (via son op.\ sur $x$) et à gauche (via
son opération sur $y$), les deux structures \struck{étant} se correspondant \uncertain{entre
elles} via la \struck{\ill{}} formule ($\ell \in \operatorname{Hom}(x_0, y_0)$,
$\gamma \in \pi$)
\[
  \ell \cdot \gamma_{x_0} = {}^{\ell}\gamma_{y_0} \cdot \ell
  \qquad \text{où } {}^{\ell}\gamma_{y_0} = \ell \cdot \gamma \cdot \ell^{-1}
\]
\note{Dans le premier membre, un symbole biffé précède le ${}^{\ell}\gamma_{y_0}$.}

----

\struck{Si $x, y \in P$, alors} L'image inverse du groupoïde $\Pi_1(P_0)$ \struck{\ill{}} par
l'application \add{$P \to P_0$} \struck{\ill{}} sur l'ens.\ des objets \uncertain{est}
$\Pi_1(P) \wedge_\pi \mathcal{E}$. Donc il est clair que l'image inverse de $\Pi_1(P_0)$ par
la composée
\[
  E \longrightarrow P \xrightarrow{\;q\;} P_0
\]
\struck{structure} s'identifie au \struck{$\mathcal{E}$-groupoïde \ill{}}
$C' \wedge_\pi \mathcal{E}$, \struck{et alors le} \add{que \ill{}, s'obtient \ill{}}
\struck{\ill{}} \add{\uncertain{dérivant}} : partir du groupoïde $C_0$ image inverse de
$\Pi_1(P_0)$ par l'application $E_0 \to P_0$. J'ai dit que celle-ci peut se construire de
façon purement algébrique en termes de $C'$ muni des structures de $\pi$-groupoïde et de
\add{groupe d'opérateurs} $\mu$, \struck{et de la structure d'ext.\ (13)} \add{en effet}
comme
\note{Un double trait vertical, dans la marge gauche, borde les cinq dernières lignes.}

\page{97}
$\mu$ opère sur $\mathcal{E}$ de façon compatible avec son opération (triviale ici) sur $\pi$,
il opère sur $C' \wedge_\pi \mathcal{E}$ par transport de structure, en opérant librement sur
\[
  \operatorname{Ob}(C' \wedge_\pi \mathcal{E}) = \operatorname{Ob} C' \overset{\mathrm{def}}{=} E' ,
\]
donc on voit que sur le \add{\uncertain{bien}} quotient $\varepsilon = \varepsilon'/\mu$, il y a
une structure de \struck{$\pi$}\add{$\mathcal{E}$}-groupoïde — plus de doute, c'est $C_0$ !
\note{Le premier membre de la formule centrée est écrit sur un mot biffé ; après « def » un
premier $E'$ est biffé. Dans « $\varepsilon = \varepsilon'/\mu$ » la lettre est petite, sans
qu'on puisse dire si c'est le $E$ des objets ou un $\varepsilon$.}

Inversement, partons d'un \add{$\mathcal{E}$-}groupoïde
\[
  C_0 \quad \text{\struck{$\longrightarrow$}}
\]
\note{La page s'arrête là ; le reste est blanc.}

\marginal{cf.\ \uncertain{démonstration} plus haut}

\page{98}
\note{Page de notes disposées en deux colonnes : à gauche des schémas, à droite un texte numéroté ; en bas un nouveau départ, au-dessus d'une base $T$. L'ordre suivi ici est : colonne de gauche, colonne de droite, bas de page.}

\begin{tikzcd}[column sep=small, row sep=small]
 & P \arrow[d, no head] & U \arrow[d] \\
E_0 \arrow[r] & P/\mu = P_0 \arrow[d] & U_0 = U/\mu \\
 & T &
\end{tikzcd}

\note{À côté de $P_0$ un mot biffé ; sous la flèche verticale de gauche, le $T$ est écrit sur une lettre biffée.}
\[
  1 \longrightarrow \mathcal{E} \longrightarrow G_0 \longrightarrow G \longrightarrow 1,
  \qquad G_0 = \pi_1(P_0), \quad G = \pi_1(T)
\]
avec $\mathcal{E} \to \mu$ et $\varphi : G_0 \to \mu$
\note{dans le schéma, une flèche verticale descend de $\mathcal{E}$ à $\mu$ et une flèche
oblique $\varphi$ va de $G_0$ à $\mu$ ; « $G_0 = \pi_1(P_0)$ » et « $G = \pi_1(T)$ » sont
écrits sous $G_0$ et $G$ avec des guillemets d'égalité.}
\[
  \text{i.e.\ on a} \quad \pi_{1,\mathrm{ab}}(C_0) \longrightarrow \mu
\]
\[
  \text{\struck{$\operatorname{Ker}\varphi$}} \qquad
  1 \longrightarrow \pi \longrightarrow \operatorname{Ker}\varphi \xrightarrow{\;\mathrm{surj}\;} G \longrightarrow 1
\]

1°) Une structure de $\mathcal{E}$-groupoïde (central) \add{$C_0$} sur $E_0$ comme ens.\
d'objets

\struck{2°)} $C$ fournit un \uncertain{prolongement} de $P_{q_0} \to T$ comme $U_0$-torseur,
et $E_0 \to P_0$ par \uncertain{morphismes} des pts par groupoïde fondamental

2°) Un \struck{fibré} $\mu$-rev.\ principal \struck{$P$} $E$ de $E_0$, défini par un système
local sur $C_0$ sur \uncertain{$P_0$}, avec op.\ de translation de $\mu$ dessus (de façon que
$C_0 \simeq C/\mu$), l'homomorphisme
\[
  \pi_1(C_0, x_0) \longrightarrow \mu \simeq \operatorname{Aut}(E_{x_0})
\]
prolongeant \add{$\uparrow \mathcal{E}$} l'homomorphisme de passage au quotient
$\mathcal{E} \to \mu$

\[
  P \xrightarrow{\;f\;} T
\]
$P$ \struck{\ill{}} torseur sous le faisceau en groupes top.\ \uncertain{localement}
\ill{} $U_T$ (syst.\ local des $\pi_1(U_t)$) est $\Pi_T$
\[
  f^*(\Pi_T) \text{ \struck{\ill{}} opère sur } \Pi_1 P, \qquad
  f^*(\Pi_T) \longrightarrow \operatorname{Int} \Pi_1 P,
  \qquad \Pi_1 P \longrightarrow \Pi_1 T
\]
$\gamma_T$ un sous-faisceau en groupes inv.\ de $U_T$ \uncertain{(central commutatif)}
\[
  U_{0T} = U_T / \gamma_T \qquad \text{syst.\ local des } \pi_1 \text{ est } \mathcal{E}_T
\]
\[
  0 \longrightarrow \Pi_T \longrightarrow \mathcal{E}_T \longrightarrow \gamma_T \longrightarrow 1
\]

\page{99}
\note{Nouveau départ, sans lien visible avec les pages 88 à 98 ; la lettre $\mathcal{T}$ est ici tracée avec un trait doublé, celle des groupoïdes de Teichmüller du titre du dossier.}

$G$ graphe MD dont tous les sommets sont de genre 0 et d'ordre $\leq 4$ (donc 3 ou 4)
\[
  S_0\mathcal{T}_G =
  \Bigl(\prod_{s \in S}
  \underbrace{S_0\mathcal{T}_{0,\hat{A}_s}}_{\substack{(\mathbf{Z}^{\hat{A}_s},\, \mu_2^{\hat{A}_s})\\ \text{-groupoïde}}}
  \Bigr)
  \wedge^{(\mathbf{Z}^{\hat{A}_s},\, \mu_2^{\hat{A}})}
  (\mathbf{Z}^{\hat{A}}, \mu_2^{\hat{A}})
\]
un $(\mathbf{Z}^{\hat{A}}, \mu_2^{\hat{A}})$-groupoïde libre, dont le quotient par
$(\mathbf{Z}^I, \mu_2^I)$ est \struck{\ill{}} $S_0\mathcal{T}_G$, et dont le quotient par
$(\mathbf{Z}^{\hat{A}}, \mu_2^{\hat{A}})$ est
\[
  R_0\mathcal{T}_G = \prod_{s \in S}
  \underbrace{R_0\mathcal{T}_{G_s^b}}_{S_0\mathcal{T}_{G_s^b}}
  = \prod_{s \in S} \mathcal{T}_{0,\hat{A}_s}
\]
(où dans le produit, on peut se borner aux $s \in S$ d'ordre 4)
\note{Le premier membre de la formule ci-dessus est écrit sur un symbole biffé ; dans le
premier $S_0\mathcal{T}_G$ de la page, l'exposant du produit $\wedge$ porte des signes
partiellement biffés.}

De façon générale, $S_0\mathcal{T}_G$ ($G$ graphe MD quelconque) est un
$(\mathbf{Z}^{\hat{A}}, \mu_2^{\hat{A}})$-groupoïde \add{$\mu_2^{\hat{A}}$-libre}
\struck{\ill{}} \struck{est un} comme \emph{$\varinjlim$} \add{(filtrante)} de
$(\mathbf{Z}^{\hat{A}}, \mu_2^{\hat{A}})$-groupoïdes \add{$\mu_2^{\hat{A}}$-}libres, si $G$ de
genre virtuel nul sur ses comp.\ connexes, si \struck{\ill{}} \add{« rigide »} la catégorie
d'indices pour le \emph{$\varinjlim$} est \ill{} — on peut sans artificialité prendre un ens.\
ordonné (celle des arbres \add{finis} \struck{réduits} — i.e.\ sans sommets d'ordre 2 — qui
sont du type $G^b$ i.e.\ \ldots{} dont l'ens.\ des bouts sur les $\pi_0(-)$ \ldots)

\page{100}
\note{Feuillet par ailleurs blanc ; deux lignes en haut à droite, d'une encre plus noire. Le dernier indice porte, après le $G$, un signe que nous ne savons pas identifier (dièse ou bécarre).}

Hom de généralisation
\[
  SP_G \longrightarrow SP_{G^{\natural}}
\]
\note{L'argument de la page 99 s'arrête en bas de page, sur « \ldots{} », et peut continuer au-delà du lot ; la page 100 n'en est pas la suite visible.}

\end{document}
